\item Luz monocromática de longitud de onda $\lambda$ incide en un par de rendijas separadas por $2.40 \times 10^{-4}\text{ m}$ y forma un patrón de interferencia en una pantalla colocada a $1.80\text{ m}$ de las rendijas. La franja brillante de primer orden está en $y = 4.52\text{ mm}$.
\begin{enumerate}
    \item Si se asume que las franjas se distribuyen de forma estrictamente lineal, estime la posición de la franja $n = 50$ multiplicando por $50.0$.
    \item Determine la tangente del ángulo $\theta$ para la primera franja brillante.
    \item Use el resultado anterior para calcular la longitud de onda exacta de la luz.
    \item Calcule el ángulo real $\theta$ para la franja de orden $50$.
    \item Determine la posición real $y$ de la franja de orden $50$ en la pantalla.
    \item Comente detalladamente la concordancia o discrepancia entre los incisos (a) y (e).
\end{enumerate}